\documentclass{article}
\usepackage{amsmath, amssymb, amsthm}
\usepackage{geometry}
\geometry{margin=1in}

\begin{document}

\title{Geometric Intuition for Lagrange Multipliers}
\author{}
\date{}
\maketitle
\Large

\section*{Key Concepts}

\subsection*{1. Gradients and Level Curves}
For a function \( f(x, y) \), the gradient \( \nabla f \) at a point is \textbf{perpendicular} to the level curve \( f(x, y) = k \) passing through that point. This is because the directional derivative in any tangent direction \( \mathbf{v} \) satisfies:
\[
\nabla f \cdot \mathbf{v} = 0.
\]

\subsection*{2. Tangency and Parallel Gradients}
When the level curve \( f(x, y) = k \) is \textbf{tangent} to the constraint \( g(x, y) = c \) at a point \( (x^*, y^*) \):
\begin{itemize}
    \item The gradients \( \nabla f \) and \( \nabla g \) must both be perpendicular to the shared tangent line.
    \item This implies \( \nabla f \parallel \nabla g \), so:
    \[
    \nabla f = \lambda \nabla g \quad \text{for some scalar } \lambda.
    \]
\end{itemize}

\subsection*{3. Example}
Maximize \( f(x, y) = x + y \) subject to \( g(x, y) = x^2 + y^2 = 1 \):
\begin{itemize}
    \item The maximum occurs where the line \( x + y = k \) is tangent to the circle \( x^2 + y^2 = 1 \).
    \item At the tangent point \( \left( \frac{\sqrt{2}}{2}, \frac{\sqrt{2}}{2} \right) \):
    \[
    \nabla f = \begin{bmatrix} 1 \\ 1 \end{bmatrix}, \quad
    \nabla g = \begin{bmatrix} 2x \\ 2y \end{bmatrix} = \begin{bmatrix} \sqrt{2} \\ \sqrt{2} \end{bmatrix}.
    \]
    \item Clearly, \( \nabla f = \lambda \nabla g \) with \( \lambda = \frac{1}{\sqrt{2}} \).
\end{itemize}

\subsection*{4. Why This Matters}
If \( \nabla f \) and \( \nabla g \) are not parallel at a point on \( g(x, y) = c \), you could move along the constraint to increase/decrease \( f \). Only when gradients are parallel (level curves tangent) is further improvement "blocked," making it a potential extremum.


\subsection*{Summary}
\[
\boxed{
\text{Tangency of level curves } \implies \nabla f \parallel \nabla g \implies \text{ Lagrange multiplier condition}
}
\]



\section*{The Gradient Alignment Condition}

From our earlier discussion, we know that at the extremum of $ f(x) $ subject to the constraint $ g(x) = c $:
\[
\nabla f = \lambda \nabla g,
\]
where $ \lambda $ is the Lagrange multiplier. This means that the gradients of $ f $ and $ g $ are parallel at the extremum.

However, this condition alone does not enforce the constraint $ g(x) = c $. To incorporate the constraint into the optimization problem, we proceed as follows.

\section*{Introducing the Constraint}

To ensure that the solution satisfies $ g(x) = c $, we introduce the term $ g(x) - c $, which measures how much the constraint is violated. At the solution, this term must equal zero:
\[
g(x) - c = 0.
\]

\section*{Constructing the Lagrangian}

We combine the objective function $ f(x) $ and the constraint $ g(x) - c $ into a single function called the \textbf{Lagrangian}:
\[
\mathcal{L}(x, \lambda) = f(x) - \lambda (g(x) - c).
\]

Here:
\begin{itemize}
    \item $ f(x) $: The original objective function to be optimized.
    \item $ -\lambda (g(x) - c) $: A penalty term that incorporates the constraint. The multiplier $ \lambda $ adjusts the weight of the constraint in the optimization.
\end{itemize}

\section*{Why Does This Work?}

The Lagrangian encodes both the gradient alignment condition and the constraint satisfaction condition. Let’s see why:

\subsection*{(a) Gradient Condition}

Take the partial derivatives of $ \mathcal{L} $ with respect to the variables $ x $:
\[
\frac{\partial \mathcal{L}}{\partial x} = \nabla f - \lambda \nabla g.
\]
Setting $ \frac{\partial \mathcal{L}}{\partial x} = 0 $ gives:
\[
\nabla f = \lambda \nabla g.
\]
This is exactly the gradient alignment condition we derived earlier!

\subsection*{(b) Constraint Satisfaction}

Take the partial derivative of $ \mathcal{L} $ with respect to $ \lambda $:
\[
\frac{\partial \mathcal{L}}{\partial \lambda} = -(g(x) - c).
\]
Setting $ \frac{\partial \mathcal{L}}{\partial \lambda} = 0 $ gives:
\[
g(x) - c = 0 \quad \text{or equivalently} \quad g(x) = c.
\]
This ensures that the constraint is satisfied.

\section*{Summary}

By constructing the Lagrangian $ \mathcal{L}(x, \lambda) = f(x) - \lambda (g(x) - c) $, we achieve two things:
\begin{enumerate}
    \item The gradient alignment condition ($ \nabla f = \lambda \nabla g $) is enforced by setting $ \frac{\partial \mathcal{L}}{\partial x} = 0 $.
    \item The constraint $ g(x) = c $ is enforced by setting $ \frac{\partial \mathcal{L}}{\partial \lambda} = 0 $.
\end{enumerate}

Thus, solving $ \nabla \mathcal{L} = 0 $ simultaneously enforces both conditions, ensuring that we find the extrema of $ f(x) $ subject to the constraint $ g(x) = c $.

\section*{Final Boxed Answer}

The Lagrangian $ \mathcal{L}(x, \lambda) = f(x) - \lambda (g(x) - c) $ is constructed so that:
\begin{enumerate}
    \item Setting $ \frac{\partial \mathcal{L}}{\partial x} = 0 $ enforces $ \nabla f = \lambda \nabla g $ (gradient alignment).
    \item Setting $ \frac{\partial \mathcal{L}}{\partial \lambda} = 0 $ enforces $ g(x) = c $ (constraint satisfaction).
\end{enumerate}

This ensures that the solution satisfies both conditions for constrained optimization:
\[
\boxed{\mathcal{L}(x, \lambda) = f(x) - \lambda (g(x) - c)}
\]


\section*{Example 1: Maximizing Volume of a Cylinder with Surface Area Constraint}

\textbf{Problem}: Find the dimensions of a cylinder (radius $ r $, height $ h $) with maximum volume, given a fixed surface area $ S = 12\pi $. \\

\textbf{Solution}:
\begin{itemize}
    \item \textbf{Objective Function}: Volume $ V = \pi r^2 h $.
    \item \textbf{Constraint}: Surface area $ 2\pi r^2 + 2\pi rh = 12\pi $.
    \item \textbf{Lagrangian}:
    \[
    \mathcal{L}(r, h, \lambda) = \pi r^2 h - \lambda (2\pi r^2 + 2\pi rh - 12\pi).
    \]
    \item \textbf{Partial Derivatives}:
    \[
    \frac{\partial \mathcal{L}}{\partial r} = 2\pi r h - \lambda (4\pi r + 2\pi h) = 0 \quad (1),
    \]
    \[
    \frac{\partial \mathcal{L}}{\partial h} = \pi r^2 - \lambda (2\pi r) = 0 \quad (2),
    \]
    \[
    \frac{\partial \mathcal{L}}{\partial \lambda} = -(2\pi r^2 + 2\pi rh - 12\pi) = 0 \quad (3).
    \]
    \item \textbf{Solve}:
    \begin{itemize}
        \item From (2): $ \pi r^2 = 2\pi r \lambda \implies r = 2\lambda $.
        \item Substitute $ \lambda = r/2 $ into (1):
        \[
        2\pi r h = \frac{r}{2} (4\pi r + 2\pi h) \implies 4h = 4r + 2h \implies h = 2r.
        \]
        \item Substitute $ h = 2r $ into (3):
        \[
        2\pi r^2 + 2\pi r(2r) = 12\pi \implies 6r^2 = 12 \implies r = \sqrt{2}, \, h = 2\sqrt{2}.
        \]
    \end{itemize}
    \item \textbf{Conclusion}: Maximum volume is achieved when $ r = \sqrt{2} $, $ h = 2\sqrt{2} $, and $ V = 4\sqrt{2}\pi $.
\end{itemize}

\section*{Example 2: Minimizing Distance from a Point to a Plane}

\textbf{Problem}: Find the point on the plane $ 2x + 3y + z = 6 $ closest to the origin. \\

\textbf{Solution}:
\begin{itemize}
    \item \textbf{Objective Function}: Distance squared $ f(x, y, z) = x^2 + y^2 + z^2 $.
    \item \textbf{Constraint}: Plane equation $ 2x + 3y + z = 6 $.
    \item \textbf{Lagrangian}:
    \[
    \mathcal{L}(x, y, z, \lambda) = x^2 + y^2 + z^2 - \lambda (2x + 3y + z - 6).
    \]
    \item \textbf{Partial Derivatives}:
    \[
    \frac{\partial \mathcal{L}}{\partial x} = 2x - 2\lambda = 0 \implies x = \lambda \quad (1),
    \]
    \[
    \frac{\partial \mathcal{L}}{\partial y} = 2y - 3\lambda = 0 \implies y = \frac{3}{2}\lambda \quad (2),
    \]
    \[
    \frac{\partial \mathcal{L}}{\partial z} = 2z - \lambda = 0 \implies z = \frac{\lambda}{2} \quad (3),
    \]
    \[
    \frac{\partial \mathcal{L}}{\partial \lambda} = -(2x + 3y + z - 6) = 0 \quad (4).
    \]
    \item \textbf{Solve}:
    \begin{itemize}
        \item Substitute (1), (2), (3) into (4):
        \[
        2\lambda + 3\left(\frac{3}{2}\lambda\right) + \frac{\lambda}{2} = 6 \implies 7\lambda = 6 \implies \lambda = \frac{6}{7}.
        \]
        \item Critical point: $ x = \frac{6}{7} $, $ y = \frac{9}{7} $, $ z = \frac{3}{7} $.
    \end{itemize}
    \item \textbf{Conclusion}: Closest point is $ \left( \frac{6}{7}, \frac{9}{7}, \frac{3}{7} \right) $, distance $ \sqrt{\frac{18}{7}} $.
\end{itemize}

\section*{Example 3: Cobb-Douglas Production Function with Budget Constraint}

\textbf{Problem}: A firm produces output $ Q = K^{1/3} L^{2/3} $, where $ K $ is capital and $ L $ is labor. Maximize $ Q $ subject to a budget constraint $ 2L + K = 12 $. \\

\textbf{Solution}:
\begin{itemize}
    \item \textbf{Objective Function}: $ Q = K^{1/3} L^{2/3} $.
    \item \textbf{Constraint}: $ 2L + K = 12 $.
    \item \textbf{Lagrangian}:
    \[
    \mathcal{L}(K, L, \lambda) = K^{1/3} L^{2/3} - \lambda (2L + K - 12).
    \]
    \item \textbf{Partial Derivatives}:
    \[
    \frac{\partial \mathcal{L}}{\partial K} = \frac{1}{3} K^{-2/3} L^{2/3} - \lambda = 0 \quad (1),
    \]
    \[
    \frac{\partial \mathcal{L}}{\partial L} = \frac{2}{3} K^{1/3} L^{-1/3} - 2\lambda = 0 \quad (2),
    \]
    \[
    \frac{\partial \mathcal{L}}{\partial \lambda} = -(2L + K - 12) = 0 \quad (3).
    \]
    \item \textbf{Solve}:
    \begin{itemize}
        \item From (1) and (2):
        \[
        \frac{1}{3} \left( \frac{L}{K} \right)^{2/3} = \lambda, \quad \frac{2}{3} \left( \frac{K}{L} \right)^{1/3} = 2\lambda.
        \]
        Divide the two equations:
        \[
        \frac{1}{2} \left( \frac{L}{K} \right) = \frac{1}{2} \implies L = K.
        \]
        \item Substitute $ L = K $ into (3):
        \[
        2K + K = 12 \implies K = 4, \, L = 4.
        \]
    \end{itemize}
    \item \textbf{Conclusion}: Maximum output $ Q = 4^{1/3} \cdot 4^{2/3} = 4 $.
\end{itemize}

\end{document}
